\chapter{Zero-b-jet signal-region optimisation}
\label{app:sr0b_optimisation}

The $\met$ threshold in $\bVetoSR$ is studied to assess whether tighter selections can improve the sensitivity at large $\beta_L^{23}$.
A higher threshold can suppress both the SM background and the destructive interference, but also reduces the retained signal.
The significance is therefore evaluated as the threshold is varied, including only the leading systematic uncertainties.

The threshold is scanned from 200 to $1500~\GeV$ in $100~\GeV$ steps.
The signal parameter points are $(\MU,\gU,\beta_L^{23})=(1.5~\TeV,0.5,1.8)$ for $\bVetoSR$-$\Res$ and $(3.0~\TeV,2.0,1.8)$ for $\bVetoSR$-$\NonRes$.
The large value of $\beta_L^{23}$ makes these benchmarks relevant to the $\bVetoSR$ regions discussed in Section~\ref{sec:ch7_sr0b}.
The scan outcomes and adopted thresholds are summarised in Table~\ref{tab:app_sr0b_scan}.

\begin{table}[htbp]
    \centering
    \small
    \caption{Signal benchmarks and preferred thresholds in the supplementary $\bVetoSR$ $\met$ scan. The last column gives the thresholds used in the final analysis.}
    \label{tab:app_sr0b_scan}
    \begin{tabular}{lccc}
        \toprule
        Region & $(\MU~[\TeV],\gU,\beta_L^{23})$ & Scan maximum & Adopted $\met$ cut \\
        \midrule
        $\bVetoSR$-$\Res$ & $(1.5,0.5,1.8)$ & $600~\GeV$ & $>200~\GeV$ \\
        $\bVetoSR$-$\NonRes$ & $(3.0,2.0,1.8)$ & $400~\GeV$ & $>400~\GeV$ \\
        \bottomrule
    \end{tabular}
\end{table}

For $\bVetoSR$-$\NonRes$, the significance is maximised at $\met>400~\GeV$.
The scan therefore supports retaining the threshold used in the final region definition.
For $\bVetoSR$-$\Res$, the maximum occurs at $\met>600~\GeV$, giving an improvement of approximately $20\%$ in the significance at the chosen benchmark.
The final resonant selection nevertheless retains $\met>200~\GeV$ to keep its kinematic requirements aligned with $\bTagSR$-$\Res$.
The higher threshold also reduces the magnitude of the destructive interference in $\bVetoSR$-$\Res$.
This demonstrates the role of $\met$ in suppressing both the SM background and the interference contribution, although the gain must be balanced against the loss of signal events.
